\section{Inductive definitions}
\label{sec:strings-02}

\begin{definition}[Inductively generated set]
Let \(U\) be a set.  Let \(B\subseteq U\) be a set of base elements, and let
\(F\) be a set of finitary operations on \(U\).  The \emph{set inductively
generated by \(B\) under \(F\)} is the intersection of all subsets \(X\) of
\(U\) such that \(B\subseteq X\) and \(X\) is closed under every operation in
\(F\).
\end{definition}

This definition is minimal: it includes the base elements and required
closures, and nothing else.  A grammar for formulas is an inductive definition
of this form.

\begin{theorem}[Structural induction]
Let \(G\) be generated from base elements \(B\) by finitary constructors
\(F\).  Let \(P(x)\) be a metalinguistic property of members of \(G\).  Suppose
that:
\begin{enumerate}
  \item \(P(b)\) holds for every \(b\in B\); and
  \item for every \(k\)-ary constructor \(f\in F\) and all
        \(x_1,\ldots,x_k\in G\), if \(P(x_i)\) holds for every
        \(1\leq i\leq k\), then \(P(f(x_1,\ldots,x_k))\) holds.
\end{enumerate}
Then \(P(x)\) holds for every \(x\in G\).
\end{theorem}
\begin{proof}
Let \(X\defeq\set{x\in G:P(x)}\).  Condition 1 gives \(B\subseteq X\), and
Condition 2 says that \(X\) is closed under every constructor in \(F\).  By the
minimality clause in the definition of \(G\), \(G\subseteq X\).  Hence
\(P(x)\) holds for every \(x\in G\).
\end{proof}

\begin{warning}
A proof by structural induction must contain one case for every base clause and
one case for every constructor clause in the grammar actually in force.  Cases
from a familiar but different grammar do not establish the result.
\end{warning}

